\documentclass[17pt]{extarticle}  % print quality
\usepackage{polyglossia}
\usepackage[margin=10mm]{geometry}
\setdefaultlanguage{marathi}
\setotherlanguages{english}
\usepackage{fourier}
%\textbf with "Sanskrit Text" produces the shape warning because the default Sanskrit Text font on Windows does not have a bold equivalent; try Yashomudra or Shobhika; now I like Tiro better
\newfontfamily\devanagarifont[Script=Devanagari]{Sanskrit Text}
%\newfontfamily\devanagarifont[Script=Devanagari]{Tiro Devanagari Marathi}

\pagenumbering{gobble}
\setlength{\parindent}{0pt}% Remove paragraph indent
\usepackage[skip=\medskipamount]{parskip}
\usepackage{fontawesome}
\usepackage[style=english]{csquotes}
\usepackage{verse}
\usepackage{xcolor}
\usepackage{hyperref}
\hypersetup{
    colorlinks=true,
    linkcolor=blue,
    filecolor=magenta,
    citecolor=blue,
    urlcolor=purple,
}

% control hyphenation: https://tex.stackexchange.com/a/177179/64425
\tolerance=1
\emergencystretch=\maxdimen
\hyphenpenalty=10000
\hbadness=10000
% control hyphenation: https://tex.stackexchange.com/a/177179/64425

\begin{document}

\title{माकडे (गझल, १९८१)}
\author{मंगेश पाडगांवकर}
\date{}
\maketitle
\hrule

थोर आपण मानिली ही भ्रष्ट सारी माकडे\\
अन् दिले हातांत त्यांच्या पेटलेले काकडे\\

कां अता छाती पिटावी लागतां आगी घरां ?\\
देव येइल कोणता हे निस्तराया साकडे?\\

धर्म यांचा मारणे सोयीप्रमाणे या उड्या\\
याचसाठी लाभलेले पाय यांना फाकडे\\

काल ज्या झाडांवरी ही आश्रयाला थांबली\\
आज ही शत्रूप्रमाणे पाहती त्यांच्याकडे\\

जा कुठेही हीच तेथे चेहरे फेंदारुनी\\
पाय जेथे टाकिती ते सर्व रस्ते वाकडे\\

पक्ष यांचा कोठला, टोळी तयांची कोठली ?\\
जात यांची एक अंती, प्रश्न सारे तोकडे\\

हालवोनी शेपट्यांना घोषणा यांनी दिल्या ...\\
आणिली तुमच्या चितेस्तव शोधुनी ही लाकडे!\\

\vspace{5mm}
\hrule
\end{document}
